\section{Implementation}
\label{sec:implementation}
%
Both improvements are local to the routine that evaluates the metric shape functions at a set of query points. Assembly, loads, boundary conditions and solvers remain unchanged. \Cref{alg:metric} summarizes the evaluation for one element. Here, the query points are passed with their reference coordinates, which are known to the element routines, e.g., as quadrature or plotting points.

\begin{algorithm}[htbp]
\caption{Metric shape functions of one element at the query points $\bxi_r$ (\PM: trial space $\Qsk$, Lagrange basis; \PH: hybrid trial space $\Tp_p$).}
\label{alg:metric}
\begin{algorithmic}[1]
\Require nodes $\bx_j$, reference nodes $\bxi_j$, query points $\bxi_r$, $\bx_r=\bx_K(\bxi_r)$
\State $\bx_0\gets\bx_K(\vect 0)$, $\matr J_0\gets\matr J(\vect 0)$; $\bxib_j\gets\matr J_0^{-\T}(\bx_j-\bx_0)$, $\bxib_r\gets\matr J_0^{-\T}(\bx_r-\bx_0)$
  \Comment{skew coordinates \cref{eq:skew}}
\State scale $\bxib_j,\bxib_r$ with the bounding box of $\{\bxib_j\}$ \Comment{\cref{eq:bbox}}
\If{\PM}
  \State $V_{cj}\gets\ell_a(\tilde\xi_j)\ell_b(\tilde\eta_j)$, $P_{cr}\gets\ell_a(\tilde\xi_r)\ell_b(\tilde\eta_r)$, $0\le a,b\le p$
    \Comment{1D GLL Lagrange polynomials}
  \State gradients: $\nabla_{\bx}=\matr J_0^{-1}\operatorname{diag}(r_{\xi}^{-1},r_{\eta}^{-1})\nabla_{\bxit}$
\Else \Comment{\PH, \cref{eq:hybrid}}
  \For{$0\le a\le p$, $0\le b\le p$}
    \If{$a+b\le p$}
      \State $V_{cj}\gets L_a(\tilde\xi_j)L_b(\tilde\eta_j)$, $P_{cr}\gets L_a(\tilde\xi_r)L_b(\tilde\eta_r)$; gradient as in line 5
    \Else
      \State $V_{cj}\gets L_a(\xi_j)L_b(\eta_j)$, $P_{cr}\gets L_a(\xi_r)L_b(\eta_r)$; $\nabla_{\bx}=\matr J(\bxi_r)^{-1}\nabla_{\bxi}$
    \EndIf
  \EndFor
\EndIf
\State stop if $\mathrm{rcond}(\matr V)<\varepsilon_{\mathrm{mach}}$ \Comment{shape functions do not exist}
\State LU factorization of $\matr V$; $\matr N_{\mathrm M}\gets\matr V^{-1}\matr P$, $\nabla\matr N_{\mathrm M}\gets\matr V^{-1}\nabla\matr P$
\end{algorithmic}
\end{algorithm}

The one-dimensional Lagrange basis costs $\mathcal O(p)$ operations per function and point (product formula, with prefix and suffix products for the derivative, so that query points may coincide with nodes), compared with $\mathcal O(1)$ for monomials or recurrences. The hybrid trial functions additionally need $\matr J(\bxi_r)$, which is computed by the element routines anyway. The dominant cost, namely the LU factorization of the $n\times n$ matrix $\matr V$ and the solves with the $\mathcal O(p^2)$ right-hand sides of $\matr P$ and its gradients, is of the order of $p^6$ for all variants. Measured run times are given in \cref{sec:num-cost}.

If $\matr V$ is singular to working precision, the element is rejected. The companion paper, in contrast, moves a corner node of such an element by a small fraction of its shortest adjacent edge~\cite{Eisentraeger2026axisymmetric}. Since this changes the mesh, we do not use it here and report rejected cases instead (\cref{sec:num-cylinder}).

For anisotropic degrees $p_\xi\ne p_\eta$, the trial space of \PM, $\Qp_{p_\xi,p_\eta}(\bar\xi,\bar\eta)$, contains the complete polynomials of degree $m=\min(p_\xi,p_\eta)$. The hybrid space is constructed in the same way with $m$ instead of $p$ in the first group of \cref{eq:hybrid} and with the complement $L_a(\xi)L_b(\eta)$, $a\le p_\xi$, $b\le p_\eta$, $a+b>m$. Hence, both trial spaces have the same complete degree $m$.
